% =========================================================
% Chapter: Methodology II — Synchronicity-Aware Multimodal Deepfake Detection
% Purpose:
% - Describe the dual-stream audio–visual architecture and how it detects cross-modal asynchrony
% - Emphasize practical, implementation-oriented descriptions (minimal math)
% - Detail preprocessing, alignment, fusion, and training decisions to ensure reproducibility
% Navigation:
% - Sections: Problem statement, data/preprocessing, architecture (streams + fusion + head),
%   biological motivation, training protocol, fairness/robustness design, and a workflow figure
% Tips:
% - Keep attention descriptions temporal/cross-modal (not spatial patch-based) unless implemented
% =========================================================
\chapter{Methodology II: Synchronicity-Aware Multimodal Deepfake Detection}
\label{ch:method-multimodal}

This chapter presents Phase II of the research: a multimodal audio–visual deepfake detector that explicitly models cross-modal temporal relationships between speech acoustics and facial articulation. The design integrates specialized per-modality encoders with bidirectional temporal modeling and fuses them via multi-head cross-modal attention to detect short-duration viseme–phoneme misalignments indicative of manipulation in contemporary multimodal deepfakes. This phase is motivated by limitations observed in the visual-only setting, where hybrid forgeries can evade detectors by preserving one modality; here we therefore model audio–visual synchrony explicitly to address such cases.

% High-level task framing: 4-class classification (RR/RF/FR/FF) enables explicit cross-modal manipulation identification
\section{Problem Formulation}
We frame multimodal deepfake detection as a four-class classification problem over manipulation categories: RR (Real-Video/Real-Audio), RF (Real-Video/Fake-Audio), FR (Fake-Video/Real-Audio), and FF (Fake-Video/Fake-Audio). Four-way classification enables explicit identification of cross-modal manipulations. For reporting a binary setting (authentic vs. manipulated), the manipulated classes (RF, FR, FF) are grouped together. Training uses a standard multi-class classification objective without presenting formal equations here.

% Data and preprocessing: define datasets, consistent sampling/alignment, and modality-specific feature preparation
\section{Data and Preprocessing}
% Benchmarks: selected for diversity in subjects and manipulation types; enables fairness and generalization assessment
\subsection{Benchmark Datasets}
To ensure diversity in manipulation types and demographics, training and evaluation employ a composite of two large-scale benchmarks:
\begin{itemize}
    \item \textbf{FakeAVCeleb:} Multimodal audio–video dataset with subjects spanning diverse races and genders, including categories RR/RF/FR/FF created using lip-sync, voice cloning, and face-swapping attacks. \cite{fakeavceleb2021}
    \item \textbf{AV-Deepfake1M++:} A stratified subset emphasizing high-fidelity, content-driven deepfakes, including advanced TTS and voice conversion attacks, as well as robust lip-sync methods. A representative sample is selected to balance feasibility and diversity.
\end{itemize}

% Visual preprocessing: standardize frame count and normalize to backbone expectations for stable training
\subsection{Video Preprocessing}
\begin{itemize}
    \item \textbf{Frame Sampling:} Uniform sampling to a fixed sequence of 16 frames.
    \item \textbf{Resolution and Normalization:} Frames resized to 224x224 and normalized to ImageNet statistics or [0,1] as required by the backbone.
\end{itemize}

% Audio preprocessing: resample, duration-normalize, and convert to Mel-spectrograms for CNN processing
\subsection{Audio Preprocessing}
\begin{itemize}
    \item \textbf{Waveform Conditioning:} Resample to 16~kHz, trim/pad to a fixed duration (e.g., 3~s).
    \item \textbf{Time–Frequency Features:} Mel-spectrogram via STFT with FFT size 2048, hop length 512, and 128 Mel bins, producing a time–frequency representation (e.g., 128 by 94). \cite{griffinlim1984}
\end{itemize}

% Augmentations: simulate real-world degradations in a modality-consistent way (apply across the sequence)
\subsection{Data Augmentation}
To improve robustness under in-the-wild degradations:
\begin{itemize}
    \item \textbf{Visual:} Random horizontal flips ($p=0.5$), small rotations ($\pm 10^\circ$), color jitter (brightness/contrast/saturation), Gaussian blur, and random erasing.
    \item \textbf{Audio:} Time stretching (rate 0.8–1.2), pitch shifting ($\pm 2$ semitones), background noise addition (SNR 20–40~dB), and SpecAugment (time/frequency masking). \cite{specaugment2019}
\end{itemize}

% Overview of the multimodal pipeline: parallel streams (video/audio), per-stream temporal modeling, cross-modal attention fusion, classification
\section{Proposed Architecture}
The framework comprises parallel modality-specific streams enhanced with bidirectional temporal modeling, fused via cross-modal attention, and followed by a multi-class head.

% Video stream: per-frame ResNet50 embeddings → Bi-GRU to capture lip dynamics and visual artifacts over time
\subsection{Video Stream}
A ResNet50 backbone processes each frame to produce frame-level spatial embeddings. In line with the paper implementation, early layers are frozen to preserve low-level features and the final stage is fine-tuned to adapt to deepfake artifacts. A two-layer Bidirectional GRU then models temporal dynamics across the sequence, capturing visual artifacts and motion patterns (e.g., blink irregularities, lip dynamics inconsistent with phoneme timing).

% Audio stream: spectrogram CNN → Bi-GRU to capture prosody and rhythm; align audio steps to video frames
\subsection{Audio Stream}
A 4-block convolutional encoder (Conv–ReLU–MaxPool) processes the Mel-spectrogram into hierarchical spectral features. A two-layer Bidirectional GRU then models prosody and rhythm over time. Audio time steps are aligned to video frame indices via consistent sampling or linear interpolation to a common sequence length.

% Cross-modal attention: each modality queries the other to localize short-duration viseme–phoneme mismatches
\subsection{Cross-Modal Attention Fusion}
To expose short-duration asynchronies, cross-modal attention allows each stream to query the other. This bi-directional attention aligns audio and video features, focusing on time spans where viseme–phoneme mismatches are likely to occur. The outputs are then pooled and concatenated to form a fused representation that emphasizes cross-modal inconsistencies \cite{katamneni2024ijcb,chen2023icassp}.

% Classification head: maps fused representation to 4 classes (RR/RF/FR/FF); binary variants derived by marginalization
\subsection{Classification Head}
A feed-forward network with dropout produces 4-class probabilities over {RR, RF, FR, FF}. Consistent with the paper implementation, the fused features are passed to a 2-layer MLP that reduces from 1536 to 256 units and then to 4 logits, followed by softmax. For binary reporting, the probabilities of manipulated classes (RF, FR, FF) are summed to obtain a single manipulated score. We avoid formal notation here and describe the head in practical terms.

% Why synchrony matters: speech co-articulation imposes timing constraints that forgeries fail to replicate uniformly
\section{Biological and Forensic Motivation}
Human speech production exhibits consistent co-articulation patterns: facial articulation (visemes) typically precedes or co-occurs with corresponding phoneme acoustics within tight temporal windows. Generative pipelines trained on separate objectives for audio and video may achieve perceptual coherence but struggle to reproduce fine-grained, biologically plausible lags and coarticulation dynamics uniformly. Modeling bidirectional temporal context with explicit cross-modal attention helps detect these subtle asynchronies, which often persist even when spatial and spectral artifacts are minimized \cite{haliassos2021lips}.

% Training details: objective, optimizer, regularization, and stability techniques; keep math minimal and implementation-focused
\section{Training Protocol and Optimization}
\begin{itemize}
    \item \textbf{Objective:} Multi-class cross-entropy over (RR, RF, FR, FF). Optionally, an auxiliary alignment loss (e.g., contrastive alignment between synchronized segments) can be explored in ablations, but is not required for the main results.
    \item \textbf{Optimization:} Adam or AdamW with initial learning rate 1e-4–3e-4, cosine decay, weight decay ~1e-4, and gradient clipping. \cite{loshchilov2019adamw}
    \item \textbf{Regularization:} Dropout in GRU and FFN layers (0.1–0.3), label smoothing (optional).
    \item \textbf{Training Details:} Mixed precision where available; per-GPU batch sizes set by memory. Early stopping on validation AUC-ROC or macro F1.
\end{itemize}

% Evaluation design: demographic stratification and compression robustness to reflect deployment realities
\section{Fairness and Robustness Evaluation Design}
\begin{itemize}
    \item \textbf{Demographic Stratification:} For FakeAVCeleb, compute metrics per subgroup (e.g., race, gender) and report subgroup accuracies, macro-averaged metrics, and variance across subgroups to assess fairness.
    \item \textbf{Compression Robustness:} Evaluate under varying compression levels and resolutions (e.g., H.264 at multiple bitrates) to approximate platform-induced degradations.
    \item \textbf{Ablations:} Compare cross-modal attention against feature concatenation; evaluate unimodal streams independently (audio-only, video-only) to quantify the multimodal gain.
\end{itemize}

% Workflow figure: end-to-end view of streams, cross-modal attention, and classification head
\begin{figure}[!htbp]
\centering
\resizebox{\linewidth}{!}{%
\begin{tikzpicture}[node distance=1.8cm]
% Video stream
\node[node] (vinput){Video Frames};
\node[node, right=of vinput] (vresnet){ResNet50\\(per-frame)};
\node[node, right=of vresnet] (vbigru){Bi-GRU\\(video)};
% Audio stream
\node[node, below=2.2cm of vresnet] (awave){Audio Waveform};
\node[node, right=of awave] (astft){STFT\\Mel-Spectrogram};
\node[node, right=of astft] (acnn){CNN\\(spectral)};
\node[node, right=of acnn] (abigru){Bi-GRU\\(audio)};
% Cross-modal attention and head
\node[node, above right=1.1cm and 1.8cm of vbigru] (xattn){Cross-Modal\\Attention (8 heads)};
\node[node, right=of xattn] (fuse){Fusion\\(pool/concat)};
\node[node, right=of fuse] (cls){Classifier\\(RR / RF / FR / FF)};

% Arrows for video path
\draw[arrow] (vinput) -- (vresnet);
\draw[arrow] (vresnet) -- (vbigru);

% Arrows for audio path
\draw[arrow] (awave) -- (astft);
\draw[arrow] (astft) -- (acnn);
\draw[arrow] (acnn) -- (abigru);

% Cross connections to attention
\draw[arrow] (vbigru) |- (xattn);
\draw[arrow] (abigru) |- (xattn);

% From attention to fusion and classifier
\draw[arrow] (xattn) -- (fuse);
\draw[arrow] (fuse) -- (cls);
\end{tikzpicture}
}%
\caption{Phase II Multimodal Architecture with Cross-Modal Attention}
\label{fig:phase2-arch}
\end{figure}

% Recap: summarizes the multimodal design and connects forward to experimental validation in Chapter~\ref{ch:experiments}
\section{Summary}
The proposed synchronicity-aware framework integrates modality-specific encoders and bidirectional temporal modeling with cross-modal attention to detect fine-grained audio–visual asynchronies characteristic of modern multimodal deepfakes. The architecture is trained with multi-class supervision to distinguish cross-modal manipulation categories and is evaluated for accuracy, robustness under compression, and fairness across demographic subgroups. Experimental results are presented in Chapter~\ref{ch:experiments} and compared against unimodal and naive fusion baselines.